% This must be in the first 5 lines to tell arXiv to use pdfLaTeX, which is strongly recommended.
\pdfoutput=1
% In particular, the hyperref package requires pdfLaTeX in order to break URLs across lines.

\documentclass[11pt]{article}

% Remove the "review" option to generate the final version.
\usepackage[]{ACL2023}

% Standard package includes
\usepackage{times}
\usepackage{latexsym}
\usepackage{amsmath}

% For proper rendering and hyphenation of words containing Latin characters (including in bib files)
\usepackage[T1]{fontenc}
% For Vietnamese characters
% \usepackage[T5]{fontenc}
% See https://www.latex-project.org/help/documentation/encguide.pdf for other character sets

% This assumes your files are encoded as UTF8
\usepackage[utf8]{inputenc}

% This is not strictly necessary, and may be commented out.
% However, it will improve the layout of the manuscript,
% and will typically save some space.
\usepackage{microtype}

% This is also not strictly necessary, and may be commented out.
% However, it will improve the aesthetics of text in
% the typewriter font.
\usepackage{inconsolata}
\usepackage{graphicx}
\usepackage{booktabs}


% If the title and author information does not fit in the area allocated, uncomment the following
%
%\setlength\titlebox{<dim>}
%
% and set <dim> to something 5cm or larger.

\title{ChatDBGPro: GDB Is All You Need}

% Author information can be set in various styles:
% For several authors from the same institution:
% \author{Author 1 \and ... \and Author n \\
%         Address line \\ ... \\ Address line}
% if the names do not fit well on one line use
%         Author 1 \\ {\bf Author 2} \\ ... \\ {\bf Author n} \\
% For authors from different institutions:
% \author{Author 1 \\ Address line \\  ... \\ Address line
%         \And  ... \And
%         Author n \\ Address line \\ ... \\ Address line}
% To start a seperate ``row'' of authors use \AND, as in
% \author{Author 1 \\ Address line \\  ... \\ Address line
%         \AND
%         Author 2 \\ Address line \\ ... \\ Address line \And
%         Author 3 \\ Address line \\ ... \\ Address line}

\author{Authors: \text{Ibraheem Amin, Shreyas Garimella, Anika Mehrotra} \\ Contact: \text{\{ia8920, sg6104, am9487\}} @princeton.edu}

\begin{document}
\maketitle
\thispagestyle{plain}
\pagestyle{plain}
\begin{abstract}
Large language models are increasingly used as debugging assistants, but it remains unclear how much of their debugging capability is attributable to model scale versus information access (tool calling and context retrieval). We extend ChatDBG, a debugger interface that pairs an LLM with GDB/LLDB tool calls, by evaluating whether smaller, cost-efficient models can match frontier performance when given appropriate tool access. We evaluate eight models -- GPT-5.5, GPT-4o, Claude Sonnet 4.5, Grok-4, Gemini 3.1 Flash-Lite, Nemotron-3-Nano-30B, Qwen3-30B, and Llama-3.1-8B -- across 20 synthetic crash cases and 19--20 real-world C/C++ bugs (BugsCPP and injected-repo), scored on three independent axes: root cause, local fix, and global fix. We compare two tool surfaces: a bash-only baseline (T1) and ChatDBG's gdb-only debugger interface (T2). We find that model scale is not the primary determinant of debugging success: with the ChatDBG tool surface, Nemotron-3-Nano-30B (3B active params, $\sim$100$\times$ cheaper than frontier) reaches 0.95/0.80/0.90 on synthetic T2, the top score in our matrix. Several frontier models regress from T1 to T2, which we attribute to bash-vs-debugger training-distribution mismatch; a structured GDB cheatsheet partially closes this gap. On real-world bugs, debugger access lifts every capable model to a 0.47--0.74 root-cause band, while Llama-3.1-8B remains at 0.00 despite engaging with tools.


\end{abstract}

\section{Methods}
\label{sec:methods}

The original ChatDBG result of 67\% actionable-fix rate \citep{levin2024chatdbg} fixes a single configuration: GPT-4 driving a sanitized gdb interface across nine single-file C/C++ crashes. Two of the three relevant variables, the model and the tool surface, are not independently controlled, so the result does not separate model capability from interface design. We address this with a factorial sweep over three axes: the buggy program ($\sim$40 cases split across synthetic and real-world corpora), the tool surface (two tiers in the headline matrix, T1 bash-only and T2 gdb-only), and the model (eight under test), evaluated end-to-end by an automated three-axis rubric. The sweep yields $\sim$40 $\times$ 2 $\times$ 8 $\approx$ 640 graded runs, plus targeted ablations on context length, gdb sandbox policy, and prompt enrichment.

\subsection{Test Corpus}
\label{sec:methods:corpus}

Cases must crash deterministically, expose the bug to runtime inspection, and have a known ground-truth fix; correctness or performance bugs that do not crash are out of scope. All cases are C or C++, both for direct comparability with ChatDBG's original evaluation and because manual memory management is the regime that exercises an interactive debugger.

The corpus splits into a 20-case synthetic set and a 19--20-case real-world set in the headline matrix, with additional cases used in ablations and extended sweeps. The synthetic set comprises nine direct ports of ChatDBG's original suite (\texttt{bench/cases/paper/}) and eleven hand-authored crashes covering bug classes underrepresented in the original (signed/unsigned loops, integer-overflow allocation, double-free on error paths, vector-iterator invalidation, off-by-one, use-after-free on linked lists, uninitialized stack reads, stack-buffer overflows). The real-world set comprises four BugBench programs \citep{lu2005bugbench}; five injected-repo cases derived by reverting known patches in pinned versions of cJSON, Lua, Mongoose, the SQLite shell, and zlib, in the spirit of the repository-level reversion methodology of SWE-bench \citep{jimenez2024swebench}; five BugsCPP libtiff cases and five BugsCPP berry cases \citep{an2023bugscpp}. Extended sweeps draw additionally on eleven external-benchmark cases (six from CrashBench \citep{ortega2024crashbench} and five from the Juliet test suite \citep{juliet2017} covering CWE-121, 122, 126, 415, and 416).

This expansion is motivated by two considerations. First, the original nine-case suite has been on GitHub since 2023 and is plausibly in the pretraining mix of every model evaluated; one memorized fix shifts a nine-case score by 11 points. Second, with five models and four tiers, a nine-case corpus produces noise wider than the effects of interest. Hand-authored cases were written in 2026 and remain unpublished, postdating every model's training cutoff. The injected-repo cases reuse bugs whose public patches predate cutoffs, but their failing input, crash signature, and exact repository state are novel; pattern recognition is plausible while verbatim retrieval is not.

Each case is described by a \texttt{case.yaml} specifying language, source, build flags, run arguments, expected-crash flag, ground-truth bug metadata, and a \texttt{criteria} block carrying three pre-registered scoring rubrics (\S\ref{sec:methods:rubric}). The criteria are visible to the judge but not to the model under test.

\subsection{Tool-Surface Tiers}
\label{sec:methods:tiers}

A \emph{tier} fixes a tool surface; the model, prompt scaffolding, context window, and sandbox image are held constant across tiers, so any tier-vs-tier difference is attributable to the interactive surface alone. The headline matrix reports two implemented tiers, T1 and T2. Two additional tiers (T3, T4) were designed to extend the comparison but were not implemented in time for this paper; we describe their intended purpose at the end of this section and revisit them as next steps.

\paragraph{Tier 1 (bash only).} A \texttt{mini-swe-agent} harness \citep{minisweagent2025} exposing a single \texttt{bash} tool. The model has no debugger and must reproduce or instrument via the shell. T1 is the no-debugger baseline and isolates how far a model can get with general-purpose shell tools alone.

\paragraph{Tier 2 (gdb only, ChatDBG native).} ChatDBG's original tool surface: a sanitized \texttt{debug} tool, \texttt{get\_code\_surrounding}, and \texttt{find\_definition}; no shell. This is the closest replication of the configuration evaluated in the original ChatDBG paper, and T2 minus T1 isolates the marginal effect of swapping a generic shell for a structured debugger surface. As an ablation we run \emph{T2-unfenced}, in which \texttt{CHATDBG\_UNSAFE=1} lifts the gdb command allow-list (\texttt{break}, \texttt{run}, \texttt{continue}, \texttt{disassemble}, \texttt{x}, \texttt{set}, \texttt{call}, full \texttt{backtrace}), measuring the cost of ChatDBG's safety sanitizer.

\paragraph{Tier 3 (bash + persistent gdb, designed).} A \texttt{mini-swe-agent} harness extending T1 with an additional \texttt{gdb} tool dispatched to a long-lived gdb subprocess (\texttt{LocalGdbBashEnvironment}), so that debugger state persists across calls and lives alongside an unrestricted shell. T3 was designed to isolate whether bash and gdb compose additively or whether the structured-only surface (T2) is on net more useful than a richer combined surface. We did not have time to implement or test T3 for this paper.

\paragraph{Tier 4 (Claude Code CLI, designed).} The Claude Code CLI run with \texttt{--bare}, providing bash, file read, and file edit. T4 was designed as a frontier-agent ceiling reference rather than a controlled comparison: strictly more tools than the other tiers, useful for bounding how far the best off-the-shelf agent gets when given full repository access. We did not have time to implement or test T4 for this paper.

Both implemented tiers run inside Docker containers (\texttt{synthetic-runner} for synthetic cases, \texttt{gdb-<project>} for BugsCPP and injected-repo cases). Patch git history is stripped at workspace prep to prevent agents from recovering the developer fix via \texttt{git log}, a leak path identified during pilot runs. Tier configurations are declarative JSON files under \texttt{bench/configs/}.

\subsection{Models}
\label{sec:methods:models}

To test whether frontier scale is necessary for tool-mediated debugging, the model lineup spans the cost-and-scale spectrum rather than concentrating at the frontier. We evaluate eight models: four frontier closed models, one frontier-small closed model, two open-weights $\sim$30B mixture-of-experts models, and one small open-weights dense model (Table~\ref{tab:models}). All models are accessed through LiteLLM (typically via OpenRouter) at temperature 0 (or the provider minimum), with a tier-calibrated max-tokens cap, so prompt and response handling is uniform across providers. GPT-4o serves a dual role: it appears in the model lineup and is also used as the judge (\S\ref{sec:methods:judge}). To avoid self-preference, when grading GPT-4o's responses the judge follows the same strict-mode 0/1 protocol used for every other model, and we hand-grade a portion of GPT-4o responses as a calibration check.

\begin{table}[t]
    \centering
    \footnotesize
    \setlength{\tabcolsep}{4pt}
    \begin{tabular}{@{}llll@{}}
        \toprule
        Model & Role & Params & Cutoff \\
        \midrule
        \texttt{gpt-5.5} & frontier & --- & 2026 \\
        \texttt{gpt-4o} & frontier (also judge) & --- & 2023 \\
        \texttt{claude-sonnet-4.5} & frontier & --- & 2025 \\
        \texttt{grok-4} & frontier & --- & 2025 \\
        \texttt{gemini-3.1-fl} & frontier-sm & --- & 2025 \\
        \texttt{nemotron-3-nano} & open MoE & 3B/30B & 2025 \\
        \texttt{qwen3-30b-a3b} & open MoE & 3B/30B & 2025 \\
        \texttt{llama-3.1-8b} & open dense & 8B & 2023 \\
        \bottomrule
    \end{tabular}
    \caption{Models under evaluation (providers: OpenAI, Anthropic, xAI, Google, NVIDIA, Qwen, Meta). All cutoffs predate or coincide with the hand-authored cases in \S\ref{sec:methods:corpus}.}
    \label{tab:models}
\end{table}

\subsection{Run Procedure}
\label{sec:methods:loop}

A run is identified by the tuple $(\text{case}, \text{tier}, \text{model}, \text{context\_lines}, \text{seed})$, hashed into a deterministic \texttt{run\_id} so re-execution is a cache hit and partially completed sweeps resume safely. Each run compiles the case in a clean Docker workspace (with bug patch applied and git history stripped, for injected-repo cases), runs the binary to crash to capture sanitizer output and backtrace, then enters the agent loop until a final response is emitted or a step, cost, or wall-clock budget is exceeded.

The system prompt (\texttt{instructions/default.txt}) provides generic debugging instructions and descriptions of the active tier's tools; the bug category, ground-truth lines, and scoring criteria are withheld. The user prompt (\texttt{build\_initial\_prompt}) supplies six fields: case metadata, an enriched stack trace (top three frames with $\pm$\texttt{context\_lines} of source, default 10), the sanitizer or crash message, the invocation command line and stdin, the debugger command history, and a fixed question, \emph{``What is the root cause of this crash? Walk through the program state, identify the defect, and propose a fix in code.''} The enriched stack trace is delivered identically on every tier, so tier-vs-tier differences isolate the interactive surface and the \texttt{context\_lines} ablation isolates static context.

\subsection{Scoring Rubric}
\label{sec:methods:rubric}

The original paper grades a single binary axis: did the model identify the root cause? This collapses three distinct outcomes (understanding, symptom-level patching, and structural fix) and obscures the failure mode most relevant to the small-model debugging hypothesis, in which a model localizes correctly but proposes the wrong kind of fix. We score every response on three independent 0/1 axes:

\begin{enumerate}
    \item \textbf{Root cause}: the response correctly diagnoses the defect (right line, variable, and invariant).
    \item \textbf{Local fix}: the proposed change resolves the proximate symptom at the crash site.
    \item \textbf{Global fix}: the proposed change addresses the bug class structurally (API, type, or invariant), not just the crashing line.
\end{enumerate}

The axes are independent; local without global is the most common outcome, and global without local is rare. The criteria block is pre-registered per case and visible only to the judge.

\subsection{LLM Judge}
\label{sec:methods:judge}

Manual grading does not scale to $\sim$600 responses, and pattern-matching is too brittle (line numbers shift; correct rationales use synonyms). We use an LLM judge, following the methodology of \citet{zheng2023judging}, which characterizes the three biases relevant here: position, verbosity, and self-preference.

The judge (\texttt{bench/judge.py}) loads each run's \texttt{case.yaml}, \texttt{result.json}, and buggy source, then issues a strict-mode prompt (\texttt{prompts/judge\_system.txt}) requiring 0/1 scores with no partial credit and a two-sentence rationale per axis. The source is truncated at 20{,}000 characters and the response at 40{,}000. The judge returns a single JSON object \texttt{\{root\_cause, local\_fix, global\_fix, rationale\}}, written to a sibling \texttt{score.json} with judge token counts for cost auditing.

We use GPT-4o as the judge for three reasons. Its October 2023 training cutoff predates the hand-authored cases (\S\ref{sec:methods:corpus}) and most evaluated model releases, limiting recency-driven self-preference. It follows JSON-schema constraints reliably (parse-failure rates near zero in pilot). Its biases are documented in the calibration literature \citep{zheng2023judging,gu2024llmjudgesurvey}, so they can be partially corrected. Bias mitigations follow from the taxonomy: the judge sees only the model's final response, never the debugger transcript, removing transcript-length bias; the strict 0/1 protocol with bounded rationales discourages hedging; and we hand-grade $\sim$10\% of responses with a target judge--human Cohen's $\kappa \geq 0.7$ before relying on the headline numbers. GPT-4o also appears in the model-under-test lineup, which introduces direct self-preference risk; for that subset specifically, we cross-checked a sample of GPT-4o responses against a stricter judge (\texttt{gpt-5}) and found the ranking preserved with GPT-4o granting itself roughly +0.05 absolute score on average, well within the spread between models in the headline matrix.


\section*{Results}
\section{Results}
\label{sec:results}

We organize results around four research questions. All scores are mean per-axis values on the 0--1 rubric (\S\ref{sec:judge}) unless stated otherwise. Parenthetical fractions denote completed (non-timeout) runs out of total attempted. All primary results use \textbf{unconditional scoring}: runs that time out or fail to produce output are scored 0/0/0, preventing survivorship bias. Ablation comparisons use conditional scores (ok runs only) and are noted as such. All primary data comes from a single curated dataset (\texttt{final\_paper\_bench}), which contains one run per (case, model, tier) cell with full provenance tracking. 

\subsection{RQ1: Model Scale vs.\ Tool Surface}
\label{sec:rq1}

Figure~\ref{fig:heatmap} shows per-axis scores across eight models, two corpora (20 synthetic, 19--20 real-world), and two tiers (T1: bash only; T2: ChatDBG/LLDB).

\paragraph{Small open-source models rival frontier with debugger tools.}
On synthetic T1, model scale dominates. Sonnet-4.5 scores 1.00/1.00/0.80 (RC/LF/GF) and GPT-5.5 scores 0.85/0.80/0.85. Nemotron-30B collapses to 0.10/0.05/0.05 (5/20 completed), effectively non-functional without debugger access. Moving to T2, the ranking shifts dramatically. Nemotron-30B (3B active parameters) achieves 0.95/0.80/0.90, the highest synthetic T2 score of any model, completing 19/20 runs. Qwen-30B (also 3B active) scores 0.85/0.65/0.65. Both are served via OpenRouter at \$0.05/M input tokens, roughly 100x cheaper than frontier alternatives. The improvement is attributable entirely to the structured ChatDBG prompt and tool surface, since the model weights are identical between T1 and T2.

\paragraph{Frontier models regress from T1 to T2 on synthetic bugs.}
Counter to the expectation that more tools should help, several frontier models score lower with ChatDBG than with bash alone on synthetic cases. Sonnet-4.5 drops from 1.00/1.00/0.80 (T1) to 0.40/0.30/0.20 (T2), a 0.63 average decline. GPT-5.5 drops from 0.85/0.80/0.85 to 0.75/0.70/0.75. Grok-4 drops from 0.75/0.65/0.70 to 0.15/0.15/0.10.

We attribute this to training distribution mismatch. Frontier models are extensively trained on bash interactions (\texttt{cat}, \texttt{grep}, \texttt{find}) but encounter GDB/LLDB commands far less frequently in their training data. In T1, they leverage deeply familiar tools. In T2, they must operate an unfamiliar interface. Grok-4 averages only 1.9 tool calls per T2 run, barely engaging with the debugger. Sonnet-4.5 makes 15.9 calls/run but its top commands (\texttt{print}=74, \texttt{code}=73) produce a score of only 0.40/0.30/0.20. It issues commands but cannot synthesize results into correct diagnoses. The GDB cheatsheet ablation (\S\ref{sec:cheatsheet}) confirms this: adding a command reference to the prompt improves GPT-4o by +0.22 (conditional), showing that syntax unfamiliarity is a binding constraint. This is not a limitation of GDB as a tool. It is a limitation of how frontier models have been trained. Small models with less bash-specific training adapt more readily to the structured debugger interface.

\paragraph{Real-world bugs favor different models.}
On real-world T1, GPT-5.5 dominates at 0.95/0.95/0.95 (19/19), while GPT-4o drops to 0.25/0.25/0.15 (20/20), a sharp reversal from its synthetic T1 performance. On real-world T2, GPT-4o recovers to 0.74/0.74/0.47 (19/19), the second-best score. Nemotron-30B achieves 0.53/0.53/0.47 on real-world T2, a large improvement from 0.05/0.05/0.05 on T1, though it no longer matches frontier as it does on synthetic cases. Llama-8B scores 0.00/0.00/0.00 on real-world T2 despite completing all 19 runs. It makes tool calls (1.9/run) but cannot interpret results on multi-file codebases.

\begin{figure*}[t]
\centering
\includegraphics[width=\textwidth]{figures/heatmap.png}
\caption{Per-axis debugging scores across models, tiers, and corpora (unconditional). Scores are mean RC/LF/GF over all attempted cases per cell; timeout or failure counts as 0. Data: \texttt{final\_paper\_bench}, May 2026.}
\label{fig:heatmap}
\end{figure*}

\subsection{RQ2: Tool Usage and Investigation Strategy}
\label{sec:rq2}

We extract per-run tool-call counts, token usage, and LLDB command distributions from the T2 transcripts in \texttt{final\_paper\_bench}.

\paragraph{More tool calls do not predict higher scores.}
On synthetic T2, Nemotron-30B achieves 0.95/0.80/0.90 with only 4.2 tool calls per run. GPT-4o averages 5.3 calls and scores 0.85/0.80/0.40. Gemini-FL-Lite makes 31.0 calls per run, 7x more, but scores only 0.60/0.55/0.45. The relationship between tool volume and score is weak and non-monotonic. What matters is which commands the model issues, not how many.

\paragraph{Three distinct strategies emerge.}
Transcript analysis reveals three behavioral clusters. The first is investigate-then-diagnose, used by GPT-5.5 (16.6 calls/run) and Sonnet-4.5 (15.9 calls/run). These models use high call volume balanced across \texttt{code}, \texttt{print}, and \texttt{bt}. GPT-5.5 reads source (100 \texttt{code} calls across 19 runs), prints variables (79), and checks backtraces (36). On real-world bugs, GPT-5.5 escalates to 78.2 calls/run and 822 \texttt{code} calls. The second is inspect-and-answer, used by Nemotron-30B (4.2 calls/run) and GPT-4o (5.3 calls/run). These models focus on the highest-value commands: Nemotron-30B issues mostly \texttt{code} (57 total) and \texttt{bt} (12), reading source around the crash site and diagnosing without exploratory frame-walking. The third is unproductive engagement, seen in Gemini-FL-Lite (31.0 calls/run with 49 \texttt{run}, 37 \texttt{break}, and 45 \texttt{print} commands, repeatedly re-running the program) and Grok-4 (1.9 calls/run, barely engaging at all).

\paragraph{Source reading dominates; execution commands do not help.}
Across all synthetic T2 runs, \texttt{code} (read source around a location) is the most-called command at 3.6/run, increasing to 12.6/run on real-world bugs. Execution-oriented commands like \texttt{run} (0.6/run), \texttt{break} (0.4/run), and \texttt{continue} (0.3/run) are used primarily by low-scoring models. The winning strategy is to inspect, not re-execute: navigate frames, read surrounding code, print variables. On real-world T2, models also resort to \texttt{shell} commands (2.0/run aggregate, led by Sonnet-4.5 at 136 calls) to work around the T2 tool constraint.

\begin{figure}[t]
\centering
\includegraphics[width=\columnwidth]{figures/fig4_gdb_commands.png}
\caption{GDB/LLDB command distribution per run, synthetic (left) and real-world (right). Source reading (\texttt{code}) dominates both corpora. Execution commands (\texttt{run}, \texttt{break}) are used primarily by low-scoring models.}
\label{fig:gdb_commands}
\end{figure}

\begin{figure}[t]
\centering
\includegraphics[width=\columnwidth]{figures/fig5_tool_calls.png}
\caption{Mean tool calls per run by model. Nemotron-30B (4.2 calls) outscores Gemini-FL-Lite (31.0 calls) on synthetic T2.}
\label{fig:tool_calls}
\end{figure}

\begin{figure}[t]
\centering
\includegraphics[width=\columnwidth]{figures/fig6_tokens.png}
\caption{Mean token usage per run by model. Sonnet-4.5 uses the most tokens (18K on real-world) but does not achieve the highest scores.}
\label{fig:tokens}
\end{figure}

\subsection{RQ3: Synthetic vs.\ Real-World Bugs}
\label{sec:rq3}

Every model's score drops from synthetic to real-world on at least one axis, but the magnitude varies with both model scale and tool mode.

\paragraph{GPT-5.5 is the only model that improves on real-world T1.}
GPT-5.5 scores 0.95/0.95/0.95 on real-world T1 versus 0.85/0.80/0.85 on synthetic T1. This may reflect its broader pretraining on real codebases. All other models degrade. GPT-4o drops from 0.80/0.75/0.35 to 0.25/0.25/0.15. Nemotron-30B drops from 0.10/0.05/0.05 to 0.05/0.05/0.05.

\paragraph{Real-world T2 compresses the model ranking.}
On real-world T2, the spread between models narrows compared to synthetic T2. Five models cluster between 0.47 and 0.74 on RC, while only Llama-8B scores 0.00. The debugger provides a floor of capability on real-world bugs that distinguishes models that can use tools at all from those that cannot.

\paragraph{Global fix remains the universal bottleneck.}
Across every condition, GF lags behind RC and LF. On synthetic T2, GPT-4o scores 0.85/0.80/0.40. It locates bugs and suggests correct fixes but explains root causes only 40\% of the time. On real-world T2, GPT-4o scores 0.74/0.74/0.47. This confirms that GF is an explanation-quality problem: articulating why a bug exists requires deeper reasoning than locating or patching it.

\begin{figure*}[t]
\centering
\includegraphics[width=\textwidth]{figures/fig2_per_axis.png}
\caption{Per-axis scores by model across all four conditions (unconditional). GF consistently lags RC and LF. Nemotron-30B shows the largest T1 to T2 improvement on synthetic cases.}
\label{fig:per_axis}
\end{figure*}

\subsection{RQ4: Ablations}
\label{sec:rq4}

We evaluate three targeted ablations on the synthetic corpus under T2. All ablation data was collected on 2026-05-12. Ablation runs used different case sets than \texttt{final\_paper\_bench}, so we report conditional scores (ok runs only) to isolate the ablation effect from completion-rate differences.

\subsubsection{GDB Cheatsheet Prompt Enrichment}
\label{sec:cheatsheet}

We augment the system prompt with a structured GDB/LLDB command reference listing common inspection commands with usage examples. This is a pure prompt intervention: no tools are added or removed. Figure~\ref{fig:cheatsheet} shows the effect.

GPT-4o shows the largest improvement (+0.22 from 0.60 to 0.82 conditional average), confirming that debugger syntax unfamiliarity constrains its T2 performance. Nemotron-30B with cheatsheet reaches 0.96, approaching frontier-level performance at a fraction of the cost. Grok-4 worsens (0.50 to 0.18): models that do not engage with tools are not helped by knowing which tools exist. The cheatsheet produces improvements comparable to switching model families, reinforcing that prompt quality is a first-order lever for debugging performance.

\begin{figure}[t]
\centering
\includegraphics[width=\columnwidth]{figures/fig8_cheatsheet.png}
\caption{GDB cheatsheet effect on mean conditional score (synthetic T2). Nemotron-30B with cheatsheet (0.96) approaches GPT-5.5 (1.00). Data: \texttt{synth-lldb-cheatsheet}, 2026-05-12.}
\label{fig:cheatsheet}
\end{figure}

\subsubsection{Context-Window Ablation}
\label{sec:context}

We vary the source lines included around each stack frame: $k \in \{5, 10, 20, 999\}$, where $k = 999$ effectively includes the entire file. The $k=10$ baseline uses \texttt{final\_paper\_bench} synthetic T2; other values use dedicated ablation sweeps. Figure~\ref{fig:context} shows the results.

Context size has non-monotonic effects. Several models score higher at $k=5$ than $k=10$, suggesting that minimal context forces more active tool use. $k=20$ is optimal for most models, providing enough local context to see the surrounding function without overwhelming the prompt. $k=999$ (full file) does not help consistently: it floods the prompt with irrelevant code, diluting the signal from the crash site. Nemotron-30B is robust across all sizes (0.84 to 0.88), consistent with its reliance on tool calls rather than prompt context.

\begin{figure}[t]
\centering
\includegraphics[width=\columnwidth]{figures/fig7_context.png}
\caption{Context-line ablation (synthetic T2). $k=20$ is the sweet spot for most models. Small models like Nemotron-30B are robust to context variation. Data: \texttt{synth-ctx\{5,20,999\}-v2}, 2026-05-12; $k=10$ from \texttt{final\_paper\_bench}.}
\label{fig:context}
\end{figure}

\subsubsection{Check-My-Work (CMW) Iterative Feedback}
\label{sec:cmw}

After the model's diagnosis, a GPT-4o judge scores RC/LF/GF and returns per-axis hints. The model revises until it achieves 3/3 or stalls for two consecutive rounds. Sample sizes are small (5--9 per model). Figure~\ref{fig:cmw} shows per-axis results.

Mid-tier models benefit most from iterative feedback. GPT-4o's GF jumps from 0.42 to 0.86 (+0.44). GPT-5.5 reaches 1.00 across all axes, typically in one to two iterations; feedback is redundant for the strongest model. Nemotron-30B declines from 0.95 to 0.67 on RC and from 0.95 to 0.56 on GF, suggesting it overwrites correct initial answers with worse revisions during the feedback loop. Llama-8B shows marginal improvement (GF: 0.06 to 0.25) but remains weak. These results confirm that GF is the axis where CMW provides the most value. It is an explanation-quality problem amenable to iterative refinement for models with sufficient reasoning capacity, but small models lack the ability to incorporate feedback productively.

\begin{figure}[h!]
\centering
\includegraphics[width=\columnwidth]{figures/fig9_cmw.png}
\caption{Check-My-Work iterative feedback effect on each scoring axis (synthetic T2, conditional). Mid-tier models (GPT-4o) benefit most on GF. Small models (Nemotron-30B) decline. Data: \texttt{synth-cmw-v2}, 2026-05-12.}
\label{fig:cmw}
\end{figure}

\subsection{Summary of Findings}
\label{sec:summary}

\begin{enumerate}
    \item \textbf{Structured tools compensate for model scale.} Nemotron-30B (3B active, \$0.05/M tokens) improves from 0.07 to 0.88 mean score (T1 to T2 synthetic) and achieves the highest synthetic T2 score of any model (0.95/0.80/0.90), exceeding GPT-5.5 (0.75/0.70/0.75).

    \item \textbf{Frontier models are trained on bash, not GDB.} Sonnet-4.5 drops 0.63 average from T1 to T2 on synthetic cases. GPT-5.5 drops 0.10. These models are fluent with shell tools but unfamiliar with debugger commands. The GDB cheatsheet partially closes this gap (+0.22 for GPT-4o).

    \item \textbf{Tool-call strategy matters more than volume.} Nemotron-30B achieves the top synthetic T2 score with 4.2 calls/run. Gemini-FL-Lite scores lower with 31 calls/run. Source reading (\texttt{code}) and backtrace (\texttt{bt}) predict success; execution commands (\texttt{run}, \texttt{break}) do not.

    \item \textbf{Global fix is the universal bottleneck.} GF lags RC and LF by 10 to 45 points across all conditions. Mid-tier models benefit most from CMW feedback on this axis (+0.44 for GPT-4o). Small models cannot self-correct.

    \item \textbf{Real-world bugs compress rankings but expose Llama-8B.} Five models cluster at 0.47 to 0.74 RC on real-world T2. Only Llama-8B scores 0.00 despite completing all runs.

    \item \textbf{Prompt quality is a first-order lever.} The GDB cheatsheet (+0.22 for GPT-4o, +0.15 for Nemotron-30B) produces improvements comparable to switching model families. Context $k=20$ outperforms both $k=10$ and $k=999$.
\end{enumerate}

\section*{Limitations}
\paragraph{Sample size and stability.} Each (model, tier, corpus) cell is the mean of
  19--20 cases under a single trial at temperature 0 (or provider minimum). Even at
  temperature 0 we observe non-deterministic responses from several providers; ablation
  sample sizes (5--9 per cell for CheckMyWork)
\section*{Next Steps}

A natural extension of this work is to replace prose-only scoring with executable patch verification. Our rubric currently grades the model's final response as text, without compiling and re-running an extracted patch. Strict patch execution would close the gap between articulated correctness and verified correctness, at the cost of penalizing prose explanations that are correct but not in a directly-applicable format. Building an extract-and-apply pipeline (patch parser, build harness, regression-test runner per case) would let us report a verified-fix axis alongside RC/LF/GF and would tighten the comparison with the original ChatDBG result, which is itself prose-graded.

\section*{Acknowledgements}
We acknowledge the use of generative AI tools for writing code, debugging code, and assisting in writing this paper. 

% Entries for the entire Anthology, followed by custom entries
\bibliography{ref}
\bibliographystyle{acl_natbib}

\end{document}
